\chapter{Relations familiales et commémorants}\label{ch:relations}

Les termes de parenté et de relation désignent, selon les textes, la personne commémorée ou celle qui commémore. Le relevé ci-dessous ne les distingue pas automatiquement : il fournit pour chaque terme l’effectif des textes de l’ensemble C, le nombre de textes où le terme est lu sans restitution, la répartition par période de ces derniers et la liste des notices (tableau~\ref{tab:relations}).

% généré par scripts/04_generated.py
\begin{table}[tbp]\centering\small\caption{Termes de parenté et de relation dans l’ensemble C (textes, non occurrences) : attestations restitutions comprises, attestations lues, et répartition des attestations lues par période.}\label{tab:relations}
\begin{tabular}{@{}l r r r r r r@{}}\toprule Terme & Total & Lu & 340–359 & 360–374 & 375–389 & 390–403\\\midrule
\lat{coniux, compar} & 39 & 22 & 5 & 4 & 7 & 6 \\
\lat{maritus} & 27 & 19 & 5 & 5 & 6 & 3 \\
\lat{uxor} & 12 & 8 & 3 & 4 & 0 & 1 \\
\lat{filius, filia} & 33 & 27 & 3 & 9 & 9 & 6 \\
\lat{pater} & 9 & 8 & 2 & 1 & 5 & 0 \\
\lat{mater} & 8 & 8 & 0 & 3 & 1 & 4 \\
\lat{parentes} & 12 & 8 & 2 & 3 & 1 & 2 \\
\lat{frater} & 6 & 4 & 0 & 1 & 1 & 2 \\
\lat{soror} & 4 & 3 & 1 & 0 & 1 & 1 \\
\lat{nepos, neptis} & 1 & 0 & 0 & 0 & 0 & 0 \\
\lat{amicus, amica} & 7 & 6 & 1 & 1 & 3 & 1 \\
\lat{socrus, gener, nurus} & 3 & 3 & 1 & 0 & 0 & 2 \\
\lat{virginius, virginia} & 5 & 3 & 1 & 0 & 2 & 0 \\
\bottomrule\end{tabular}\end{table}


\section{Le couple au centre de la commémoration}

Les termes conjugaux (\emph{coniux}, \emph{compar}, \emph{maritus}, \emph{uxor}, \emph{virginius}) dominent nettement les autres relations. La commémoration d’un conjoint par l’autre est la configuration la plus fréquente parmi les textes qui nomment un commanditaire : Agrippina pour Felix (\nnum{133}), Theophila pour Benenatus (\nnum{110}), Constantia pour Anastasius (\nnum{43}, en partie restitué), Victora pour Iunianus (\nnum{4}), Aradius Melissus pour Esuvia Ianuaria (\nnum{26}), Sabbatius pour Staturina (\nnum{386}), Viricunda pour Severianus (\nnum{400}). Plusieurs textes associent le commanditaire à la sépulture : Marcus fait pour Karisia et pour lui-même une « demeure éternelle » (\emph{feci nobis domum aeternalem}, \nnum{46}) ; le prêtre Gaudentius grave l’épitaphe de son épouse et la sienne (\nnum{310}) ; Clodius Adelfius fait de même (\nnum{36}).

La part des textes qui mentionnent un terme conjugal décroît au cours de la période (chapitre~\ref{ch:formulaire}). Il n’est pas possible de dire si ce recul traduit un changement des pratiques de commémoration, une modification de la composition sociale ou topographique du corpus, ou l’état de conservation des textes les plus tardifs.

\section{Parents et enfants}

Les parents commémorent leurs enfants, le plus souvent sans être nommés (\emph{parentes fecerunt}, \emph{parentes filio dulcissimo}) : Cheretius et Severa pour leur fils Cheretius, mort à cinq ans (\nnum{17}) ; les parents de Successa, morte à deux ans (\nnum{35}) ; les « malheureux parents » du \nnum{93}. Les termes \emph{mater} et \emph{pater} désignent tantôt la personne commémorée (Eleuteria, pour qui son fils Pacatianus fait un emplacement, \nnum{73} ; Nunita, mère de Maximilla, \nnum{306}), tantôt le commanditaire (le père Eutyches, \nnum{68} ; le père Nicetius, \nnum{296} ; « la mère et le père » de Felicitas, \nnum{470}). L’épitaphe de Proiecta (\nnum{244}), due à l’évêque Damase, nomme le mari, Primus, la mère et le père, Florus.

Le \nnum{399} réunit, sur une même plaque, Caelestina, épouse d’Ampliatus, morte à vingt ans, et Hilaritas, \emph{filia} morte à vingt-trois ans, qui ne peut donc être la fille de Caelestina ; toutes deux étaient mariées, et elles sont mortes à quelques semaines d’intervalle ; le \nnum{134}, un prêtre, un diacre et leur sœur. Ces textes à plusieurs défunts sont des documents familiaux autant que funéraires.

\section{Au-delà de la famille}

L’amitié est rarement invoquée, mais elle peut fonder une sépulture. Lucceia ensevelit dans son propre sarcophage la vierge Maximilla, \emph{propter amicitiam sibi coniunctam} (\nnum{306}). Cinnamius Opas, lecteur, et les époux du \nnum{4} sont dits « amis des pauvres » (\emph{amicus}, \emph{amator}, \emph{amatrix pauperum}) ; Quirillus, « aimé de tous » (\emph{omnibus amicatus}, \nnum{275}). Dans les textes où ils sont lus, les termes \emph{frater} et \emph{soror} se comprennent naturellement comme des liens de parenté (Priscus pour sa sœur, \nnum{132} ; deux frères qui se préparent une sépulture de leur vivant, \nnum{441}) ; rien n’y indique un emploi ecclésial.

% généré par scripts/04_generated.py
\begin{small}\noindent \lat{coniux, compar} : 1, 10, 13, 19, 26, 28, 33, 40, 43, 46, 47, 49, 83, 115, 117, 125, 130, 147, 184, 200, 242, 262, 275, 288, 304, 306, 310, 319, 320, 324, 333, 364, 365, 378, 382, 399, 420, 455, 472.\par\noindent \lat{maritus} : 21, 26, 29, 32, 42, 53, 77, 82, 114, 125, 133, 144, 184, 185, 193, 219, 223, 244, 254, 257, 291, 297, 324, 373, 386, 399, 454.\par\noindent \lat{uxor} : 10, 23, 34, 36, 43, 56, 104, 110, 146, 167, 400, 438.\par\noindent \lat{filius, filia} : 6, 13, 35, 42, 59, 64, 68, 81, 84, 92, 99, 100, 127, 143, 186, 196, 211, 213, 221, 233, 259, 268, 297, 306, 309, 347, 353, 378, 399, 402, 412, 431, 457.\par\noindent \lat{pater} : 2, 22, 68, 203, 241, 268, 296, 297, 378.\par\noindent \lat{mater} : 73, 81, 102, 306, 353, 382, 461, 470.\par\noindent \lat{parentes} : 17, 35, 64, 74, 93, 124, 126, 152, 189, 244, 409, 439.\par\noindent \lat{frater} : 132, 247, 314, 368, 441, 457.\par\noindent \lat{soror} : 12, 134, 181, 469.\par\noindent \lat{nepos, neptis} : 378.\par\noindent \lat{amicus, amica} : 22, 133, 189, 202, 214, 268, 391.\par\noindent \lat{socrus, gener, nurus} : 48, 378, 382.\par\noindent \lat{virginius, virginia} : 4, 45, 218, 265, 286.\end{small}

